\section{IOP with Generalized Oracles}
\label{sec:oracle}

\subsection{Generalized Oracles}
\label{sec:oracleSub}

An oracle is characterized by its message space $\mathbb{M}$, query space $\mathbb{Q}$, and the answer function $\mathcal{F} : \mathbb{M} \times \mathbb{Q} \rightarrow \Field$.
We focus on settings where $\mathbb{M}$ consists of vectors or matrices, where vectors are a special case of matrices.
The following oracle types are commonly studied:

\begin{itemize}
    \item \emph{Point oracle}:  
    $\MSpace = \Field^{m \times n}$, $\QSpace = \mathbb{N} \times \mathbb{N}$,  
    $\AFunc(\bm{M}, (i,j)) = m_{ij}$,  
    where $\bm{M} = \{m_{ij}\}_{i=1,j=1}^{m,n} \in \MSpace$ and $(i,j) \in \QSpace$.

    \item \emph{MLE oracle}:  
    $\MSpace = \Field^{m \times n}$, $\QSpace = \TensorSpace(m) \times \TensorSpace(n)$,  
    $\AFunc(\bm{M}, (\bm{l}, \bm{r})) = \bm{l}^\top \bm{M} \bm{r}$.

    \item \emph{Linear oracle}:  
    $\MSpace = \Field^{m \times n}$, $\QSpace = \Field^{m \times n}$,  
    $\AFunc(\bm{M}, \bm{Q}) = \langle \bm{M}, \bm{Q} \rangle$.
\end{itemize}

In this section, we consider a more general answer function,
corresponding to matrix-function oracles, defined as follows:

\begin{definition}[Matrix-Function Oracle]
\label{def:matrixOracle}
A \emph{matrix-function oracle (MFO)} is an oracle specified by the tuple
$(\{m_i, n_i\}_{i=1}^I, \{m_j', n_j'\}_{j=1}^J, f)$, where:
\begin{itemize}
    \item The message space is $\MSpace \coloneqq \bigotimes_{i=1}^I \Field^{m_i \times n_i}$, containing $I$ matrices;
    \item The query space is $\QSpace \coloneqq \bigotimes_{j=1}^J \Field^{m_j' \times n_j'}$, containing $J$ matrices;
    \item The answer function is defined as:
    \[
    \AFunc\big((\bm{M}_1, \ldots, \bm{M}_I), (\bm{Q}_1, \ldots, \bm{Q}_J)\big) \coloneqq f(\bm{M}_1, \ldots, \bm{M}_I, \bm{Q}_1, \ldots, \bm{Q}_J),
    \]
    where $f$ is a matrix function composed of operations including addition, subtraction, scalar multiplication, transposition, matrix multiplication, inner product, Hadamard product, tensor product, concatenation, and truncation.
\end{itemize}
\end{definition}

In the following,
we denote the point oracle by $\Oracle_\mathsf{P}$,
the MLE oracle by $\Oracle_\mathsf{MLE}$,
the linear oracle by $\Oracle_\mathsf{LIN}$,
and the matrix-function oracle by $\Oracle_\mathsf{MFO}$.
Constructions of such oracles have been studied in the literature.

\begin{lemma}[Point Oracle from Merkle Tree~\cite{crypto/Merkle89}]
\label{lemma:oraclePoint}
A Merkle tree, together with Merkle path authentication proofs, realizes a point oracle with computational soundness error $\kappa_\mathsf{H}$, where $\kappa_\mathsf{H}$ is the probability that a PPT adversary finds either a collision or a preimage in the underlying hash function $\mathsf{H}$.
\end{lemma}

\begin{lemma}[MLE Oracle from a Linear Code]
\label{lemma:oracleMLE}
Let $\Code : \Field^k \rightarrow \Field^\mu$ be a linear code with relative distance $\delta$,
and let a proximity parameter $\epsilon$ satisfy $0<\epsilon < \delta/2$.
There exists an MLE oracle
with negligible soundness error and knowledge-soundness error
for the following proximity relation:
    \[
        \Relation_\mathsf{prox} =
        \left\{
        \Code,
        \Root, \bm{M}~:~
        \mathsf{ptree}
        \mid
        \Root = \mathsf{ptree}.\mathsf{root}
        \wedge
        \delta_H(\mathsf{ptree}.\mathsf{leaf}, \Enc[\Code](\bm{M})) \le \epsilon
        \right\}.
    \]
Here, $\mathsf{ptree}$ denotes a partial Merkle tree with $\mu$ leaves, where some nodes may be
$\mathsf{null}$ and the remaining non-null nodes form a subtree.
The Hamming distance $\Delta_H$ is extended to vectors in $(\Field \cup \mathsf{null})^\mu$ in the natural way.
\end{lemma}

Such MLE oracles are typically instantiated from an interactive oracle proof of proximity (IOPP) for the following MLE-IOPP relation:
\begin{equation}
\label{eq:relationMLEiopp}
    \Relation_\mathsf{MLE\mbox{-}IOPP}[\epsilon]
    =
    \RelationExp{
    \begin{aligned}
        & \Code, \bm{C}, \epsilon\\
        & v, \bm{l}, \bm{r}
    \end{aligned}
    }{\bm{M}}{
    \begin{aligned}
        & v =\bm{l}^\top \bm{M} \bm{r} \\
        & \wedge~ \delta_H(\bm{C}, \Enc[\Code](\bm{M})) \le \epsilon
    \end{aligned}
    }.
\end{equation}

\begin{lemma}[IOPP to MLE Oracle]
\label{lemma:MLEfromIOPP}
Let $\Code : \Field^k \rightarrow \Field^\mu$ be a linear code with relative distance $\delta$, and suppose there exists an MLE-IOPP for $\Relation_\mathsf{MLE\mbox{-}IOPP}[\epsilon]$ for some $0 < \epsilon < \delta/2$.
Then, the following algorithms define an MLE oracle:
\begin{itemize}
    \item $\Oracle_\mathsf{MLE}.\Setup(1^\lambda, \MSpace, \QSpace) \rightarrow \pp$:  
    Run the setup algorithm of the IOPP to generate $\pp$.

    \item $\Oracle_\mathsf{MLE}.\Com(\pp, \bm{M}) \rightarrow \Oracle_\mathsf{P}(\bm{C})$:  
    Compute the codeword $\bm{C} \leftarrow \Enc[\Code](\bm{M})$ and commit to it via a point oracle $\Oracle_\mathsf{P}(\bm{C})$.

    \item $\Oracle_\mathsf{MLE}.\Query(\pp, (\bm{l}, \bm{r})) \rightarrow (v, \mathsf{bit})$:  
    Compute $v \leftarrow \bm{l}^\top \bm{M} \bm{r}$ and run the IOPP on input $(\Code, \Oracle_\mathsf{P}(\bm{C}), \epsilon, v, \bm{l}, \bm{r}; \bm{M})$,
    where any access to a position of $\bm{C}$ is handled by querying the point oracle $\Oracle_\mathsf{P}(\bm{C})$.
    Output $\TRUE$ iff the IOPP accepts.
\end{itemize}
\end{lemma}

This result has been widely employed in the context of hash-based MLPC constructions,
\eg,~\cite{iacr/Ben-SassonBHR18,eurocrypt/BrehmCFRRZ25,crypto/GolovnevLSTW23,crypto/XieZS22,crypto/ZeilbergerCF24}.
To make our exposition self-contained, we provide a proof of Lemma~\ref{lemma:MLEfromIOPP} in Appendix~\ref{appx:proofMLEfromIOPP}.

The code-switching technique (see Section~\ref{sec:codeSwitch}) enables the construction of polylogarithmic-size MLE-IOPPs for arbitrary linear codes $\Code$,
thereby enabling MLE oracle $\Oracle_\mathsf{MLE}[\Code]$ for arbitrary linear code.

It is also possible to construct an MLE oracle for some proximity parameter $\epsilon > \delta/2$ by leveraging out-of-domain (ODD) sampling (Lemma \ref{lemma:outofdomain}).
Intuitively, we commit to an MLE evaluation $\hat{v}$ at verifier-chosen random points $\hat{\bm{l}}$ and $\hat{\bm{r}}$ such that different witnesses within the list-decoding radius $\delta_\mathsf{LD}$ are distinguishable.
\begin{lemma}[MLE Oracle with ODD Sampling]
\label{lemma:odd}
Let $\Code : \Field^k \to \Field^\mu$ be a linear code with list-decoding 
radius $\delta_{\mathsf{LD}}$.  
Suppose there exists an IOPP for the following MLE-ODD-IOPP relation:
\begin{equation}
\label{eq:relationMLEoddiopp}
\Relation_{\mathsf{MLE\mbox{-}ODD\mbox{-}IOPP}}[\epsilon]
=
\RelationExp{
\begin{aligned}
    &\Code,\, \bm{C},\, \epsilon \\
    &\hat{v}, \hat{\bm{l}}, \hat{\bm{r}}, v, \bm{l}, \bm{r}
\end{aligned}
}{\bm{M}}{
\begin{aligned}
    &\hat{v} = \hat{\bm{l}}^{\top}\bm{M}\hat{\bm{r}}
    \;\wedge\;
    v = \bm{l}^{\top}\bm{M}\bm{r} \\
    &\wedge\;
    \delta_H(\bm{C}, \Enc[\Code](\bm{M})) \le \epsilon
\end{aligned}
}.
\end{equation}
Then the following algorithms define an MLE oracle for the proximity relation:
\[
\Relation_{\mathsf{prox}}
=
\left\{
\begin{aligned}
&\Code,\, \Root,\, \bm{M} \\
&\hat{v},\, \hat{\bm{l}},\, \hat{\bm{r}}
\end{aligned}
:\;
\mathsf{ptree}
\;\middle|\;
\begin{aligned}
&\Root = \mathsf{ptree}.\mathsf{root}
\;\wedge\;
\delta_H(\mathsf{ptree}.\mathsf{leaf}, \Enc[\Code](\bm{M})) \le \epsilon,\\
&\hat{v} = \hat{\bm{l}}^{\top}\bm{M}\hat{\bm{r}}
\end{aligned}
\right\}.
\]
\begin{itemize}
    \item $\Oracle_{\mathsf{MLE}}.\Setup(1^\lambda, \MSpace, \QSpace) \to \pp$:  
    Run the IOPP setup algorithm to generate $\pp$.

    \item $\Oracle_{\mathsf{MLE}}.\Com(\pp, \bm{M}) \to (\Oracle_{\mathsf{P}}(\bm{C}), \hat{v})$:  
    Compute $\bm{C} \leftarrow \Enc[\Code](\bm{M})$ and commit to it via a point 
    oracle $\Oracle_{\mathsf{P}}(\bm{C})$.  
    After receiving verifier challenges $(\hat{\bm{l}},\hat{\bm{r}})$, compute  
    $\hat{v} \leftarrow \hat{\bm{l}}^{\top}\bm{M}\hat{\bm{r}}$ and include 
    $\hat{v}$ in the commitment.

    \item $\Oracle_{\mathsf{MLE}}.\Query(\pp, (\bm{l},\bm{r})) \to (v,\mathsf{bit})$:  
    Compute $v \leftarrow \bm{l}^{\top}\bm{M}\bm{r}$ and run the IOPP on input  
    $(\Code, \Oracle_{\mathsf{P}}(\bm{C}), \epsilon, v, \bm{l}, \bm{r}; \bm{M})$, 
    where each access to $\bm{C}$ is answered via $\Oracle_{\mathsf{P}}(\bm{C})$.  
    Output $\TRUE$ iff the IOPP accepts.
\end{itemize}
\end{lemma}

Building upon MLE oracles, TensorSwitch~\cite{iacr/BunzFRW25} utilizes linear oracles as follows.

\begin{lemma}[Linear Oracle from an MLE Oracle]
\label{lemma:oracleLinear}
Given an MLE oracle $\Oracle_{\mathsf{MLE}}$, one can construct a linear oracle
$\Oracle_{\mathsf{LIN}}$ that uses the same $\Setup$ and $\Com$ algorithms.
The query algorithm $\Oracle_{\mathsf{LIN}}.\Query$ first executes the
inner-product protocol in Fact~\ref{fact:ip}, and then invokes
$\Oracle_{\mathsf{MLE}}.\Query$.
\end{lemma}

Generalizing linear oracles, we can use the following
matrix-function proof to construct matrix-function oracles from MLE oracles.

\begin{fact}[\textbf{Matrix-Function Proof}~\cite{asiacrypt/CCYY25}]
\label{fact:matrixproof}
For any matrix function $f$ as defined in Definition~\ref{def:matrixOracle},
there exists a proof that reduces the verification of $f$ to checking the MLE
evaluations of $\{\bm{M}_i \in \Field^{m_i\times n_i}\}_{i=1}^I$ and $\{\bm{Q}_j \in \Field^{m_j'\times n_j'}\}_{j=1}^J$.
The prover runs in
$O\!\left(\sum_{i=1}^I m_i n_i + \sum_{j=1}^J m_j' n_j'\right)$ field multiplications,
while the proof size and verifier time are both $O(\log n)$.
\end{fact}

\begin{lemma}[Matrix-Function Oracle from MLE Oracles]
\label{lemma:matrix_function_queries}
Let $\{\Oracle_{\mathsf{MLE}}^{(i)}\}_{i=1}^I$ be MLE oracles for message matrices
$\{\bm{M}_i \in \MSpace_i\}_{i=1}^I$ with proximity relations $\{\Relation_{\mathsf{prox}i}\}_{i=1}^I$.
Let
$f(\bm{M}_1,\dots,\bm{M}_I,\bm{Q}_1,\dots,\bm{Q}_J)$
be a matrix function as defined in Definition~\ref{def:matrixOracle} with query matrices
$\{\bm{Q}_j \in \QSpace_j\}_{j=1}^J$.
Then there exists a matrix function oracle $\Oracle_{\mathsf{MFO}}$ with message set $\mathbb{M} = \{\bm{M}_i \in \MSpace_i\}_{i=1}^I$, query set $\mathbb{Q} = \{\bm{Q}_j \in \QSpace_j\}_{j=1}^J$,
and proximity relation:
\[
\Relation_{\mathsf{prox}}
=
\left\{
\{\Root_i,\bm{M}_i\}_{i=1}^I \;\middle|\;
\bigwedge_{i=1}^I (\Root_i,\bm{M}_i)\in\Relation_{\mathsf{prox}i}
\right\}.
\]
The oracle $\Oracle_{\mathsf{MFO}}$ is constructed as follows:
\begin{itemize}
    \item
    $\Oracle_{\mathsf{MFO}}.\Setup(1^\lambda,\{\MSpace_i\}_{i=1}^I,\{\QSpace_j\}_{j=1}^J)$:
    Runs $\Oracle_{\mathsf{MLE}}^{(i)}.\Setup(1^\lambda,\MSpace_i)$ for each $i\in[I]$
    and outputs $\pp=\{\pp_i\}_{i=1}^I$.
    \item
    $\Oracle_{\mathsf{MFO}}.\Com(\pp,\{\bm{M}_i\}_{i=1}^I)$: Runs
    $\Oracle_{\mathsf{MLE}}^{(i)}.\Com(\pp_i,\bm{M}_i)$ for each $i\in[I]$
    and outputs $\Root=\{\Root_i\}_{i=1}^I$.
    \item
    $\Oracle_{\mathsf{MFO}}.\Query(\pp,\{\bm{Q}_j\}_{j=1}^J; \{\bm{M}_i\}_{i=1}^I)$: Computes
    $v \leftarrow f(\bm{M}_1,\dots,\bm{M}_I,\bm{Q}_1,\dots,\bm{Q}_J)$,
    runs the matrix-function proof in Fact~\ref{fact:matrixproof},
    and invokes $\Oracle_{\mathsf{MLE}}^{(i)}.\Query(\pp_i, (\bm{l}_i, \bm{r}_i); \bm{M}_i)$ for all $i\in[I]$ with $(\bm{l}_i,\bm{r}_i)$ from the matrix-function proof.
    The verifier additionally checks the MLE evaluations of
    $\{\bm{Q}_j\}_{j=1}^J$ and accepts iff all pass.
\end{itemize}
\iffalse
Let $T_{Qi}$, $S_i$, and $T_{Vi}$ denote the prover time, proof size, and verifier
time of $\Oracle_{\mathsf{MLE}}^{(i)}.\Query$, respectively, and let $T_{Ej}$
denote the cost of computing the MLE of $\bm{Q}_j$.
Then the query complexity of $\Oracle_{\mathsf{MFO}}$ is as follows:
\begin{itemize}
    \item \emph{Prover time:}
    $\sum_{i=1}^I T_{Qi}
    + \sum_{i=1}^I m_i n_i
    + \sum_{j=1}^J m_j' n_j'$ field multiplications;
    \item \emph{Proof size:}
    $\sum_{i=1}^I S_i
    + \sum_{i=1}^I \log(m_i n_i)
    + \sum_{j=1}^J \log(m_j' n_j')$ field elements;
    \item \emph{Verifier time:}
    $\sum_{i=1}^I T_{Vi}
    + \sum_{j=1}^J T_{Ej}
    + \sum_{i=1}^I \log(m_i n_i)
    + \sum_{j=1}^J \log(m_j' n_j')$ field multiplications.
\end{itemize}
\fi
\end{lemma}

The security of the resulting matrix-function oracle follows from the general security guarantees established in Lemma~\ref{lemma:csPlusKS}.
In Section~\ref{sec:iopsecurity}, we further present additional general security guarantees for IOPs.

\subsection{Security Guarantees for IOP}
\label{sec:iopsecurity}

\begin{lemma}[Computational Sound Protocol + Knowledge Sound Oracle = Knowledge Sound IOP]
\label{lemma:csPlusKS}
Let $\Pi$ be a protocol that is computationally sound with soundness error
$\kappa_\Pi^{\mathsf{s}}$ for the relation:
\[
    \Relation_\Pi =
    \RelationExp{
    \text{\sf public input}
    }{\bm{M}}{\text{\sf constraint}}.
\]
Assume that the verifier’s final check in $\Pi$ corresponds to an oracle query
to a knowledge-sound oracle $\Oracle$ for $\bm{M}$, with knowledge-soundness
error $\kappa_\Oracle^{\mathsf{ks}}$.
Then the following three-step protocol $\Pi'$ is knowledge sound with
knowledge-soundness error at most
$\kappa_\Pi^{\mathsf{s}} + \kappa_\Oracle^{\mathsf{ks}}$:
\begin{itemize}
    \item {\em Commit:} the prover sends the oracle commitment $\Oracle(\bm{M})$ to the verifier;
    \item {\em Reduce:} the prover and verifier execute $\Pi$;
    \item {\em Query:} the prover and verifier execute $\Oracle.\Query$.
\end{itemize}
\end{lemma}

\begin{proof}
\emph{Completeness.}
Completeness follows directly from the completeness of $\Pi$ and of the oracle $\Oracle$.

\smallskip
\noindent
\emph{Knowledge soundness.}
Let $\Prover^\ast$ cause the verifier to accept
with probability $\varrho$.
By the knowledge soundness of $\Oracle$, an extractor outputs a witness $\bm{M}$
with probability at least $\varrho - \kappa_\Oracle^{\mathsf{ks}}$
such that $\Pi$ accepts on input $(\text{\sf public input}, \bm{M})$.
By the computational soundness of $\Pi$, the probability that
$\Pi$ accepts while $(\text{\sf public input}, \bm{M}) \notin \Relation_\Pi$
is at most $\kappa_\Pi^{\mathsf{s}}$.
Therefore, the extractor outputs $\bm{M}$ satisfying
$(\text{\sf public input}, \bm{M}) \in \Relation_\Pi$
with probability at least
$\varrho - \kappa_\Oracle^{\mathsf{ks}} - \kappa_\Pi^{\mathsf{s}}$.
\end{proof}


\begin{lemma}[Computational/Knowledge Sound Protocol + Knowledge Sound Oracle = Computational/Knowledge Sound IOP]
\label{lemma:csksPlusKS}
Let $\Pi$ be a protocol that is computationally sound/knowledge sound with
soundness errors $\kappa_\Pi^{\mathsf{s}}$/$\kappa_\Pi^{\mathsf{ks}}$,
respectively.
Suppose that in $\Pi$ the prover sends a message $\bm{M}$, and that all subsequent
verifier checks involving $\bm{M}$ can be expressed as oracle query responses to
a knowledge-sound oracle $\Oracle$ with knowledge-soundness error
$\kappa_\Oracle^{\mathsf{ks}}$.

Consider the protocol $\Pi_{\mathsf{IOP}}$ obtained by replacing the transmission
of $\bm{M}$ with an oracle commitment $\Oracle(\bm{M})$, and replacing all
verifier checks involving $\bm{M}$ with executions of $\Oracle.\Query$.
Then $\Pi_{\mathsf{IOP}}$ is an IOP that is computationally and knowledge sound
with soundness errors
$\kappa_\Pi^{\mathsf{s}} + \kappa_\Oracle^{\mathsf{ks}}$
and
$\kappa_\Pi^{\mathsf{ks}} + \kappa_\Oracle^{\mathsf{ks}}$,
respectively.
\end{lemma}

The proof of this lemma is provided in Appendix~\ref{appx:proofCSKSplusKS}.


\begin{theorem}[Transforming Interactive Protocols into IOPs]
\label{theorem:iop}
Let $\Pi$ be an interactive protocol that is computationally and knowledge sound
with errors $\kappa_\Pi^{\mathsf{s}}$ and $\kappa_\Pi^{\mathsf{ks}}$,
respectively.
Its transcript has the form:
\(
\pi = (\alpha_0, x_1, \alpha_1, \ldots, x_s, \alpha_s),
\)
where $\{\alpha_j\}_{j=0}^s$ are prover messages, and
$\{x_j\}_{j=1}^s$ are verifier challenges.

Suppose that a subset of prover messages
$\{\alpha_j\}_{j\in\mathcal{J}}$ is replaced by oracle commitments
$\{\Oracle_j\}_{j\in\mathcal{J}}$, where each oracle $\Oracle_j$ is knowledge sound
with error $\kappa_{\Oracle_j}^{\mathsf{ks}}$,
and that all verifier checks involving such messages can be implemented via queries to the corresponding
oracle.
Let $\Pi_{\mathsf{IOP}}$ denote the resulting protocol.

Then $\Pi_{\mathsf{IOP}}$ is an IOP that is
computationally sound with error at most
\(
\kappa_\Pi^{\mathsf{s}} + \sum_{j\in\mathcal{J}} \kappa_{\Oracle_j}^{\mathsf{ks}},
\)
and knowledge sound with error at most
\(
\kappa_\Pi^{\mathsf{ks}} + \sum_{j\in\mathcal{J}} \kappa_{\Oracle_j}^{\mathsf{ks}}.
\)
\end{theorem}

Proof of this theorem is provided in Appendix~\ref{appx:proofToIOP}.

With these general security results,
we can recursively construct more and more complex IOPs from simpler proofs.
An example is the code-switching technique proposed in~\cite{jacm/RonZewiR24}.

\subsection{Code Switching via Matrix-Function Oracle Queries}
\label{sec:codeSwitch}

In this section, we present the code-switching framework of~\cite{jacm/RonZewiR24} by replacing the message transmission in an MLE-IOPP with efficient MLE oracles.
The more complex constructions in Sections~\ref{sec:interleave} and~\ref{sec:tensorswitch} then repeatedly apply the methodology developed here.


By Lemma~\ref{lemma:MLEfromIOPP}, for any linear code $\Code$, an MLE oracle can be
derived from an IOPP for the relation
$\Relation_{\mathsf{MLE\mbox{-}IOPP}}[\epsilon]$:
\[
\RelationExp{
\begin{aligned}
    & \Code,\ \bm{C}\in\Field^{\mu_1\times\mu_2},\ \epsilon\\
    & v\in\Field,\ \bm{l}\in\TensorSpace(k_1),\ \bm{r}\in\TensorSpace(k_2)
\end{aligned}
}{\bm{M}\in\Field^{k_1\times k_2}}{
\begin{aligned}
    & v=\bm{l}^\top \bm{M}\bm{r} \\
    & \wedge\ \Delta_H(\bm{C},\Enc[\Code](\bm{M}))\le\epsilon .
\end{aligned}
}
\]

The proximity condition
$\Delta_H(\bm{C},\Enc[\Code](\bm{M}))\le\epsilon$
can be verified by sampling $\ell$ random positions of
$\bm{C}$ and $\Enc[\Code](\bm{M})$.
If $\Delta_H(\bm{C},\Enc[\Code](\bm{M}))>\epsilon$, then with high probability
the two matrices differ on at least one of the sampled positions.

Let $\mathcal{I}\subseteq[\mu_1]$ and $\mathcal{J}\subseteq[\mu_2]$ be the row and
column index sets corresponding to the sampled positions.
Define selection matrices
$\bm{\eta}_L\in\{0,1\}^{\ell\times\mu_1}$ and
$\bm{\eta}_R\in\{0,1\}^{\ell\times\mu_2}$ such that the $i$-th row of
$\bm{\eta}_L$ (\resp, $\bm{\eta}_R$) contains a single $1$ at the index specified
by the $i$-th element of $\mathcal{I}$ (\resp, $\mathcal{J}$).
The selected entries of $\Enc[\Code](\bm{M})$, denoted by $\bm{u}$, can then be
written as:
\[
\bm{u}
=
\diag(\bm{\eta}_L \bm{C}\bm{\eta}_R^\top)
=
\diag(\bm{\eta}_L \Enc[\Code](\bm{M})\bm{\eta}_R^\top).
\]

Sampling a random tensor-span vector $\bm{\xi}\in\TensorSpace(\ell)$, this equality can be
batched into a single scalar check:
\[
u
=
\langle \bm{u},\bm{\xi}\rangle
=
\left\langle \diag(\bm{\eta}_L \bm{C}\bm{\eta}_R^\top),\bm{\xi}\right\rangle
=
\left\langle \diag(\bm{\eta}_L \Enc[\Code](\bm{M})\bm{\eta}_R^\top),\bm{\xi}\right\rangle .
\]

Since $\bm{C}$ is public in the $\Relation_{\mathsf{MLE\mbox{-}IOPP}}[\epsilon]$,
if the prover additionally sends $\bm{M}$ at the beginning of the protocol,
the verifier can check the above equality directly.
This yields a computationally and knowledge-sound MLE-IOPP, as formalized in
Algorithm~\ref{alg:iopp}.
Though this basic construction is not succinct, it serves as a
starting point for the more efficient MLE-IOPPs developed in subsequent
sections.

\begin{lemma}[Basic MLE-IOPP]
\label{lemma:iopp-soundness}
Let $\Code$ be a linear code and $0<\epsilon<1$ be a proximity parameter.
Algorithm~\ref{alg:iopp} is a computationally sound and knowledge-sound IOPP,
with computational and knowledge soundness errors bounded by:
\[
\kappa_{\mathsf{IOPP}}^{\mathsf{s}}
=
\kappa_{\mathsf{IOPP}}^{\mathsf{ks}}
\;\le\;
(1-\epsilon)^{\ell}
+
\frac{\lceil \log_2 \ell \rceil}{|\Field|-2}.
\]
\end{lemma}

Proof of Lemma~\ref{lemma:iopp-soundness} is provided in Appendix~\ref{appx:proofIOPPbasic}.

\input{alg/alg-iopp}

Combining Lemma~\ref{lemma:iopp-soundness} and Lemma~\ref{lemma:MLEfromIOPP}, we obtain that for any linear code,
one can construct an MLE oracle utilizing a point oracle to the codeword.
However, this construction is not necessarily succinct: in Algorithm~\ref{alg:iopp},
both the proof size and the verifier time are linear in the message size.

This non-succinctness arises solely from transmitting the message $\bm{M}$.
We observe that checking $\bm{M}$ in Line~\ref{line:iopp:checkEval} is an MLE
evaluation, while checking $\Enc[\Code](\bm{M})$ in
Line~\ref{line:iopp:checkM} is a matrix-function check, since
$\Enc[\Code](\bm{M})$ can be expressed as a matrix multiplication involving
$\bm{M}$ and the generator matrix $\bm{G}$ of $\Code$.
Hence, a natural way to restore succinctness is to replace the explicit
transmission of $\bm{M}$ and $\bm{G}$ with MLE oracles.


Suppose there exists an MLE oracle $\Oracle_{\mathsf{MLE}}[\Code'](\bm{M})$
for $\bm{M}$, based on another linear code $\Code'$, that admits succinct
query proofs and verifier time, and an MLE oracle
$\Oracle_{\mathsf{MLE}}(\bm{G})$ for the generator matrix $\bm{G}$.
Then, by Theorem~\ref{theorem:iop}, we can replace the explicit transmission
of $\bm{M}$ with a commitment to $\Oracle_{\mathsf{MLE}}[\Code'](\bm{M})$
while preserving both computational and knowledge soundness.

Blaze provide condidates for such succinct MLE oracles.
\begin{lemma}[MLE Oracle from Random Foldable Codes~\cite{crypto/ZeilbergerCF24}]
\label{lemma:oracleMLPCBaseFold}
There exists an MLE oracle based on a \emph{random foldable code (RFC)}
$\Code_{\mathsf{RFC}}$~\cite[Definition~9]{crypto/ZeilbergerCF24}
with $O(n\log n)$ prover time and
$O(\log n \cdot \lambda)$ proof size and verifier time,
and soundness and knowledge-soundness errors of $2^{-\lambda}$.
\end{lemma}

\begin{lemma}[MLE Oracle for the Generator Matrix]
\label{lemma:mleGenerator}
If a linear code $\Code: \Field^{k} \mapsto \Field^{\mu}$ admits $O(k)$-time
encoding, then its generator matrix $\bm{G} \in \Field^{\mu \times k}$
admits an MLE oracle with $O(k)$ prover time and succinct proof
size and verifier time per query.
\end{lemma}

\begin{proof}
Observe that
\(
    \bm{l}^\top \bm{G} \bm{r}
    = \langle \bm{l}, \Enc[\Code](\bm{r}) \rangle,
\)
where the right-hand side can be computed in $O(k)$ time.
Thus, an MLE oracle for $\bm{G}$ follows directly from a succinct proof of
the relation $v = \langle \bm{l}, \Enc[\Code](\bm{r}) \rangle$.
\end{proof}

\iffalse
\begin{corollary}[MLE Oracles for Generator Matrices of Interleaved and Tensor Codes]
\label{corollary:mleGeneratorInterleave}
Let $\Code_1: \Field^{k_1} \mapsto \Field^{\mu_1}$ and $\Code_2: \Field^{k_2} \mapsto \Field^{\mu_2}$ be linear codes that admit
linear-time encoding.
Then the generator matrices of the interleaved code $\Code_1^t$
and the tensor code $\Code_1 \kronecker \Code_2$ admit MLE oracles with
$O(k_1)$ and $O(k_1+k_2)$ prover time per query and succinct proof size and verifier time.
\end{corollary}

\begin{proof}
Let $\bm{G}_1$ and $\bm{G}_2$ denote the generator matrix of $\Code_1$.
The generator matrix of the interleaved code $\Code_1^t$ is
$\bm{G}_1 \kronecker \bm{I}_t$, where $\bm{I}_t \in \TensorSpace(t)$ is the $t \times t$ identity matrix,
and the generator matrix of the tensor code $\Code_1 \kronecker \Code_2$
is $\bm{G}_1 \kronecker \bm{G}_2$.
An MLE query to a Kronecker product matrix can be reduced to MLE queries
to its Kronecker components, \ie, $\bm{G}_1$ and $\bm{G}_2$. (see~\cite[Section~5.2]{asiacrypt/CCYY25}).
\end{proof}
\fi

\begin{corollary}[MLE Oracles for Generator Matrices ofTensor Codes]
\label{corollary:mleGeneratorInterleave}
Let $\Code_1: \Field^{k_1} \mapsto \Field^{\mu_1}$ and $\Code_2: \Field^{k_2} \mapsto \Field^{\mu_2}$ be linear codes that admit
linear-time encoding.
Then the generator matrices of the tensor code $\Code_1 \kronecker \Code_2$ admit MLE oracles with
$O(k_1+k_2)$ prover time per query and succinct proof size and verifier time.
\end{corollary}

\begin{proof}
Let $\bm{G}_1$ and $\bm{G}_2$ denote the generator matrix of $\Code_1$.
The generator matrix of the tensor code $\Code_1 \kronecker \Code_2$
is $\bm{G}_1 \kronecker \bm{G}_2$.
An MLE query to a Kronecker product matrix can be reduced to MLE queries
to its Kronecker components, \ie, $\bm{G}_1$ and $\bm{G}_2$. (see~\cite[Section~5.2]{asiacrypt/CCYY25}).
\end{proof}


Algorithm~\ref{alg:basicMLEOracle} is a line-by-line adaptation of Algorithm~\ref{alg:iopp}, in which Lines~\ref{line:iopp:sendM}, \ref{line:iopp:checkC}, \ref{line:iopp:checkM}, and~\ref{line:iopp:checkEval} are replaced by oracle commitments and oracle queries.
By Theorem~\ref{theorem:iop}, the resulting MLE oracle,
whose $\Query$ procedure is given by Algorithm~\ref{alg:basicMLEOracle},
preserves computational and knowledge soundness.
When $\Code'$ is instantiated as an RFC code, Algorithm~\ref{alg:basicMLEOracle} achieves polylogarithmic proof size and verifier time.

\begin{lemma}[MLE Oracle via Code Switching]
\label{lemma:codeswitch}
Let
$\Code:\Field^{k_1\times k_2}\to\Field^{\mu_1\times\mu_2}$
and
$\Code':\Field^{k_1\times k_2}\to\Field^{\mu_1'\times\mu_2'}$
be linear codes.
Suppose that an efficient MLE oracle for $\Code'$ is available, and let
$0<\epsilon<1$ be the proximity parameter.
Then Algorithm~\ref{alg:basicMLEOracle} yields an MLE oracle for $\Code$ whose
computational and knowledge soundness errors are both bounded by:
\[
(1-\epsilon)^{\ell}
+
\frac{\lceil \log_2 \ell \rceil}{|\Field|-2}
+
\kappa_{\Code'}^{\mathsf{ks}},
\]
where $\kappa_{\Code'}^{\mathsf{ks}}$ denotes the knowledge-soundness error of the
MLE oracle for $\Code'$.
\end{lemma}

\begin{proof}
This follows directly from Lemma~\ref{lemma:iopp-soundness} and
Lemma~\ref{lemma:csksPlusKS}.
\end{proof}

By instantiating $\Code$ and $\Code'$ in Lemma~\ref{lemma:codeswitch} with different
linear codes, we obtain a broad class of MLE oracles.
In \methodname, we switch from the interleaved code
$\Code_{\mathsf{base}}^{t}$ of a base code $\Code_{\mathsf{base}}$
to the tensor code $\Code^{\kronecker 2}$ of a second code $\Code$,
and then further switch to the interleaved RFC code
$\Code'^{(\ell+1)}$ of a third linear code $\Code'$
that admits efficient MLE queries (\eg, an RFC code).
